\documentclass[aspectratio=169,10pt]{beamer}
\usetheme[numbering=fraction,progressbar=frametitle]{metropolis}
\usepackage{fontspec}
\setsansfont{Latin Modern Sans}
\usepackage{amsmath,amssymb,mathtools,graphicx,booktabs}
\usepackage{polyglossia}
\setdefaultlanguage{english}
\setbeamertemplate{navigation symbols}{}
\setbeamerfont{frametitle}{size=\large}
\setbeamerfont{normal text}{size=\normalsize}
\setlength{\parskip}{.45em}
\hypersetup{pdfauthor={},pdfsubject={Lecture slides},pdfkeywords={},pdftitle={Complex numbers: the geometric point of view}}
\providecommand{\N}{\mathbb{N}}
\providecommand{\Z}{\mathbb{Z}}
\providecommand{\D}{\mathbb{D}}
\providecommand{\Q}{\mathbb{Q}}
\providecommand{\R}{\mathbb{R}}
\providecommand{\C}{\mathbb{C}}
\providecommand{\vect}[1]{\overrightarrow{#1}}
\providecommand{\vi}{\vec{\imath}}
\providecommand{\vj}{\vec{\jmath}}
\providecommand{\vk}{\vec{k}}
\providecommand{\norm}[1]{\left\lVert #1\right\rVert}
\providecommand{\vc}[2]{\begin{pmatrix}#1\\#2\end{pmatrix}}
\providecommand{\deter}[2]{\det\left(#1,#2\right)}
\providecommand{\abs}[1]{\left\lvert #1\right\rvert}
\providecommand{\dd}{\,\mathrm{d}}
\providecommand{\e}{\mathrm{e}}
\providecommand{\Card}{\operatorname{Card}}
\providecommand{\E}{\operatorname{E}}
\providecommand{\V}{\operatorname{V}}
\providecommand{\degre}{^{\circ}}
\providecommand{\un}[1]{\,\mathrm{#1}}
\providecommand{\sci}[2]{#1\times10^{#2}}
\providecommand{\pgcd}{\operatorname{PGCD}}
\providecommand{\ppcm}{\operatorname{PPCM}}
\title{Complex numbers: the geometric point of view}
\subtitle{Lecture and worked examples}
\date{}
\author{}
\begin{document}
\begin{frame}[plain]
\vfill
{\usebeamerfont{title}\inserttitle\par}
\medskip
{\usebeamerfont{subtitle}\insertsubtitle\par}
\medskip
{\usebeamercolor[fg]{progress bar}\rule{\textwidth}{.5pt}}
\vfill
\end{frame}
\begin{frame}{Objectives}
Represent a complex number by a point or a vector; determine a modulus and an argument; move from the algebraic form to the trigonometric form and back; interpret a modulus as a distance.
\end{frame}
\begin{frame}{Introductory activity}
The plane has an orthonormal frame \((O,\vect u,\vect v)\). To a complex number \(x+iy\) we associate the point with coordinates \((x,y)\).
\end{frame}
\begin{frame}{Introductory activity (2)}
\begin{enumerate}
\setcounter{enumi}{0}\item Plot the points \(A\), \(B\), \(C\) associated with \(a=2+i\), \(b=i\,a\) and \(c=(1+i)\,a\). Give \(b\) and \(c\) in algebraic form.

\item Compute the distances \(OA\), \(OB\) and \(OC\).

\item Compute the scalar product \(\vect{OA}\cdot\vect{OB}\). What can be said about the lines \((OA)\) and \((OB)\)?

\item Compute \(\cos\left(\vect{OA},\vect{OC}\right)\) with the scalar product, then the sign of the determinant \(\deter{\vect{OA}}{\vect{OC}}\). Deduce a measure of the oriented angle \(\left(\vect{OA},\vect{OC}\right)\).

\end{enumerate}
\end{frame}
\begin{frame}{Introductory activity (3)}
\begin{enumerate}
\setcounter{enumi}{4}\item Start again with \(a=-1+3i\) (or a point of your choice). State a conjecture about the geometric effect of a multiplication by \(i\), then by \(1+i\).

\item Which quantities attached to a complex number would have to be defined to state this conjecture in general?

\end{enumerate}
\end{frame}
\begin{frame}{Complex plane}
The plane has a direct orthonormal frame \((O,\vect u,\vect v)\). To the complex number \(z=a+ib\) (\(a\), \(b\) real) we associate the point \(M(a,b)\), called the \emph{image} of \(z\), and the vector \(\vect w\vc ab\). We say that \(z\) is the \emph{affix} of \(M\) and of \(\vect w\). The plane with this correspondence is called the \emph{complex plane}.
\end{frame}
\begin{frame}{Complex plane}
If \(A\) and \(B\) have affixes \(z_A\) and \(z_B\), the vector \(\vect{AB}\) has affix \(z_B-z_A\) and the midpoint of \([AB]\) has affix \(\dfrac{z_A+z_B}2\). The affix of \(\vect w+\vect{w'}\) is the sum of the affixes, that of \(k\vect w\) (\(k\) real) is \(k\) times the affix of \(\vect w\).

The image of \(\overline z\) is the reflection of the image of \(z\) in the horizontal axis; the image of \(-z\) is its reflection in \(O\).
\end{frame}
\begin{frame}{Modulus}
The \emph{modulus} of \(z=a+ib\) is the nonnegative real number \(\abs z=\sqrt{a^2+b^2}\). If \(M\) is the image of \(z\), \(\abs z=OM\). For a real number, the modulus is the absolute value.
\end{frame}
\begin{frame}{Modulus}
For every complex number \(z\), \(\abs z^2=z\overline z\), and \(\abs{\overline z}=\abs{-z}=\abs z\). If \(A\) and \(B\) have affixes \(z_A\) and \(z_B\), then \(AB=\abs{z_B-z_A}\).
\end{frame}
\begin{frame}{Modulus}
With \(z=a+ib\), \(z\overline z=a^2+b^2=\abs z^2\) (chapter 1). The distance \(AB\) is the norm of the vector \(\vect{AB}\), whose affix is \(z_B-z_A\).
\end{frame}
\begin{frame}{Modulus}
For all complex numbers \(z\), \(z'\) and every natural number \(n\):
\[\abs{zz'}=\abs z\abs{z'},\qquad\abs{z^n}=\abs z^n,\qquad\text{and, if }z'\neq0,\quad\abs{\dfrac1{z'}}=\dfrac1{\abs{z'}},\quad\abs{\dfrac z{z'}}=\dfrac{\abs z}{\abs{z'}}.\]
\end{frame}
\begin{frame}{Modulus}
\emph{Product.} \(\abs{zz'}^2=zz'\,\overline{zz'}=z\overline z\,z'\overline{z'}=\abs z^2\abs{z'}^2\). Moduli are nonnegative, so \(\abs{zz'}=\abs z\abs{z'}\).

\emph{Power.} By induction: \(\abs{z^0}=1=\abs z^0\), and if \(\abs{z^n}=\abs z^n\), then \(\abs{z^{n+1}}=\abs{z^n}\abs z=\abs z^{n+1}\).

\emph{Inverse and quotient.} \(\abs{z'}\times\abs{\dfrac1{z'}}=\abs{z'\times\dfrac1{z'}}=1\), then \(\abs{\dfrac z{z'}}=\abs z\abs{\dfrac1{z'}}\).
\end{frame}
\begin{frame}{Modulus}
For all complex numbers \(z\) and \(z'\), \(\abs{z+z'}\le\abs z+\abs{z'}\). The problem of the chapter gives a proof and the case of equality.
\end{frame}
\begin{frame}{Modulus}
We write \(\mathbb U\) for the set of complex numbers of modulus \(1\). Their images form the circle with centre \(O\) and radius \(1\).
\end{frame}
\begin{frame}{Modulus}
If \(z\) and \(z'\) belong to \(\mathbb U\), then \(zz'\in\mathbb U\) and \(\dfrac1z\in\mathbb U\); moreover \(\dfrac1z=\overline z\).
\end{frame}
\begin{frame}{Modulus}
\(\abs{zz'}=1\times1=1\) and \(\abs{\dfrac1z}=\dfrac11=1\). Finally \(z\overline z=\abs z^2=1\), so \(\overline z=\dfrac1z\).
\end{frame}
\begin{frame}{Modulus}
Let \(A\) and \(B\) have affixes \(a\) and \(b\), and \(M\) have affix \(z\). The equality \(\abs{z-a}=r\) (\(r>0\)) means \(AM=r\): \(M\) describes the circle with centre \(A\) and radius \(r\). The equality \(\abs{z-a}=\abs{z-b}\) (\(a\neq b\)) means \(AM=BM\): \(M\) describes the perpendicular bisector of \([AB]\).
\end{frame}
\begin{frame}{Argument and trigonometric form}
Let \(z\) be a nonzero complex number with image \(M\). An \emph{argument} of \(z\) is any measure, in radians, of the oriented angle \(\left(\vect u,\vect{OM}\right)\). We write \(\arg z\) for one of them; the others differ from it by a multiple of \(2\pi\), which we write \(\arg z\equiv\theta\pmod{2\pi}\). The number \(0\) has no argument.
\end{frame}
\begin{frame}{Argument and trigonometric form}
Every nonzero complex number \(z=a+ib\) can be written \(z=r(\cos\theta+i\sin\theta)\) with \(r=\abs z>0\) and \(\theta\equiv\arg z\pmod{2\pi}\); then
\[\cos\theta=\dfrac ar,\qquad\sin\theta=\dfrac br.\]
Conversely, if \(z=r(\cos\theta+i\sin\theta)\) with \(r>0\), then \(\abs z=r\) and \(\theta\) is an argument of \(z\). Two trigonometric forms \(r(\cos\theta+i\sin\theta)\) and \(r'(\cos\theta'+i\sin\theta')\) are equal if and only if \(r=r'\) and \(\theta\equiv\theta'\pmod{2\pi}\).
\end{frame}
\begin{frame}{Argument and trigonometric form}
For \(z\neq0\): \(\arg\overline z\equiv-\arg z\) and \(\arg(-z)\equiv\arg z+\pi\pmod{2\pi}\). A positive real number has argument \(0\), a negative real number \(\pi\); \(i\) has argument \(\dfrac\pi2\).
\end{frame}
\begin{frame}{Argument and trigonometric form}
Compute \(r=\abs z\), then \(\dfrac ar\) and \(\dfrac br\), and recognise the cosine and sine of a standard angle.
\end{frame}
\begin{frame}{Argument and trigonometric form}
\(1+i\sqrt3\): \(r=2\), \(\cos\theta=\dfrac12\), \(\sin\theta=\dfrac{\sqrt3}2\), so \(1+i\sqrt3=2\left(\cos\dfrac\pi3+i\sin\dfrac\pi3\right)\).

\(-1+i=\sqrt2\left(\cos\dfrac{3\pi}4+i\sin\dfrac{3\pi}4\right)\) ; \(-3=3(\cos\pi+i\sin\pi)\) ; \(-2i=2\left(\cos\left(-\dfrac\pi2\right)+i\sin\left(-\dfrac\pi2\right)\right)\).
\end{frame}
\begin{frame}{Argument and trigonometric form}
\begin{columns}[T,onlytextwidth]
\begin{column}{.57\textwidth}
\(-2\left(\cos\dfrac\pi5+i\sin\dfrac\pi5\right)\) is not a trigonometric form, because \(-2<0\). We write \(-2\left(\cos\dfrac\pi5+i\sin\dfrac\pi5\right)=2\left(\cos\dfrac{6\pi}5+i\sin\dfrac{6\pi}5\right)\).
\end{column}
\begin{column}{.4\textwidth}
\centering
\includegraphics[width=\linewidth,height=4.5cm,keepaspectratio]{figures/expertes-02-1-en.pdf}
\end{column}
\end{columns}
\end{frame}
\begin{frame}{Argument and trigonometric form}
\(\abs z=OM\), \(\abs{z_B-z_A}=AB\), \(\abs z^2=z\overline z\); the modulus of a product is the product of the moduli; \(z=r(\cos\theta+i\sin\theta)\) with \(r>0\); an argument is only defined up to a multiple of \(2\pi\).
\end{frame}
\end{document}
